% sec:schw-coordinates — Coordinate Systems and the Coordinate Singularity
\section{Coordinate Systems and the Coordinate Singularity}
\label{sec:schw-coordinates}

The Schwarzschild metric (\cref{eq:schwarzschild-metric}) has a
problem at $r = 2M$: the component $g_{rr} = (1 - 2M/r)^{-1}$
diverges, and $g_{tt} = -(1 - 2M/r)$ vanishes.  A geodesic integrator
that uses these coordinates will grind to a halt as a photon approaches
the horizon --- step sizes shrink to zero, errors accumulate, and the
integration never completes.  The previous section established that this
divergence is a coordinate artefact: the Kretschmann scalar is finite
at $r = 2M$.  This section constructs coordinates that eliminate the
artefact, giving the integrator a smooth path through the horizon.

\paragraph{The tortoise coordinate.}

Null geodesics reveal how a coordinate system encodes causal structure:
if we can find coordinates in which they look simple, the metric will
too.  Setting $\d s^2 = 0$ and $\d\theta = \d\varphi = 0$ in the
Schwarzschild metric gives, for $r \neq 2M$,
%
\begin{equation}
  \frac{\d t}{\d r} = \pm \left(1 - \frac{2M}{r}\right)^{-1} .
  \label{eq:schw-null-radial}
\end{equation}
%
Outside the horizon ($r > 2M$), the plus sign describes outgoing rays
and the minus sign ingoing ones.  As $r \to 2M^+$, the coordinate
velocity $\d r/\d t$ vanishes: light appears to freeze at the horizon
in Schwarzschild coordinates.  This isn't physical --- a freely falling
observer crosses the horizon in finite proper time --- but it reveals
that the $t$ coordinate is poorly adapted to the horizon geometry.
\physnote{The ``freezing'' of light at the horizon is an artefact of
Schwarzschild time $t$, which measures the proper time of a static
observer at infinity.  No static observer exists at $r = 2M$, so $t$
loses its physical meaning there.}

To untangle the coordinate pathology, we integrate
\cref{eq:schw-null-radial}.  Substituting $\d r_* = (1 - 2M/r)^{-1}\,
\d r$ and integrating gives the \emph{tortoise coordinate}
(for $r > 2M$):
%
\begin{equation}
  r_* = r + 2M \ln\!\left(\frac{r}{2M} - 1\right) .
  \label{eq:tortoise}
\end{equation}
%
The name reflects the coordinate's behaviour: as $r \to 2M^+$,
$r_* \to -\infty$, stretching the finite coordinate interval
$[2M, \infty)$ into $(-\infty, \infty)$.  Since $\d r_* = (1 -
2M/r)^{-1}\,\d r$, \cref{eq:schw-null-radial} becomes $\d t/\d r_*
= \pm 1$, so radial null geodesics are simply $t = \pm r_* +
\text{const}$ --- straight lines at $45^\circ$ in the $(t, r_*)$
plane, just as in flat spacetime.  The coordinate singularity has been
pushed to $r_* = -\infty$, but the metric still diverges there.

\paragraph{Ingoing Eddington--Finkelstein coordinates.}

The tortoise coordinate untangles the null structure but doesn't
eliminate the singularity.  The decisive step combines $t$ and $r_*$
into a single null coordinate that remains finite at the horizon.
Define the \emph{advanced time}
%
\begin{equation}
  v = t + r_* \,.
  \label{eq:ef-v-def}
\end{equation}
%
Along an ingoing radial null geodesic, $t = -r_* + \text{const}$, so
$v = \text{const}$: the coordinate $v$ labels ingoing null rays.
Substituting $\d t = \d v - \d r_*$ into the Schwarzschild line
element, with $\d r_* = (1 - 2M/r)^{-1}\,\d r$, gives
%
\begin{align}
  \d s^2 &= -\!\left(1 - \frac{2M}{r}\right)
    \bigl[\d v - (1 - 2M/r)^{-1}\,\d r\bigr]^2
    \notag \\
    &\qquad
    + \left(1 - \frac{2M}{r}\right)^{-1}\d r^2
    + r^2\,\d\Omega^2 \notag \\
  &= -\!\left(1 - \frac{2M}{r}\right)\d v^2
    + 2\,\d v\,\d r \notag \\
    &\qquad
    \underbrace{{}- (1 - 2M/r)^{-1}\,\d r^2
    + (1 - 2M/r)^{-1}\,\d r^2}_{= \; 0}
    + r^2\,\d\Omega^2 \,.
    \label{eq:ef-substitution}
\end{align}
%
The two $\d r^2$ terms cancel exactly --- one from expanding the
square, the other from the original $g_{rr}$ --- leaving the
cross-term $2\,\d v\,\d r$ as the sole radial contribution:
%
\begin{equation}
  \d s^2 = -\!\left(1 - \frac{2M}{r}\right)\d v^2
           + 2\,\d v\,\d r
           + r^2\!\left(\d\theta^2
             + \sin^2\!\theta\;\d\varphi^2\right) .
  \label{eq:ef-ingoing}
\end{equation}
%
Every component of the metric is finite at $r = 2M$: $g_{vv} = 0$,
$g_{vr} = 1$, and the angular part is $4M^2\,\d\Omega^2$.  The
coordinate singularity has been removed entirely.
\histnote{Eddington introduced these coordinates in 1924
\cite{Eddington1924}; Finkelstein rediscovered them in 1958 and
was the first to interpret $r = 2M$ as a one-way membrane rather than
a singularity \cite{Finkelstein1958}.}

The inverse metric is needed for raising indices and for the geodesic
Hamiltonian.  The $(v, r)$ block of $\metric$ has determinant $-1$,
and inverting gives $g^{vv} = 0$, $g^{vr} = 1$,
$g^{rr} = 1 - 2M/r$, with the standard angular components
$g^{\theta\theta} = 1/r^2$ and
$g^{\varphi\varphi} = 1/(r^2\sin^2\!\theta)$.  At the horizon,
$g^{rr} = 0$ while $g^{vr} = 1$: the surface $r = 2M$ is null, not
singular.
\warnnote{The cross-term $g_{vr}$ means that $v$ is not a time
coordinate in the usual sense --- surfaces of constant $v$ are null,
not spacelike.  This complicates the initial-value formulation but
poses no difficulty for geodesic integration.}

\paragraph{Light cone behaviour at the horizon.}

The physical character of $r = 2M$ becomes transparent in ingoing
Eddington--Finkelstein coordinates.  Setting $\d s^2 = 0$ and
$\d\Omega = 0$ in \cref{eq:ef-ingoing} gives
$\d v\bigl[2\,\d r - (1 - 2M/r)\,\d v\bigr] = 0$, which factors
into two families of radial null geodesics:
%
\begin{equation}
  \left(1 - \frac{2M}{r}\right)\d v = 2\,\d r
  \qquad \text{or} \qquad \d v = 0 \,.
  \label{eq:ef-null-geodesics}
\end{equation}
%
The first family has $\d v/\d r = 2/(1 - 2M/r)$: outgoing rays.  The
second family has $v = \text{const}$: ingoing rays.  Outside the
horizon ($r > 2M$), both families behave as expected --- outgoing rays
escape to larger $r$, ingoing rays fall to smaller $r$.  At $r = 2M$,
the outgoing family has $\d v/\d r \to \infty$: the light cone tilts
so that the ``outgoing'' direction is tangent to the horizon.  Inside
the horizon ($r < 2M$), the factor $(1 - 2M/r)$ is negative.
Future-directed null vectors satisfy $\d v > 0$ (inherited from the
orientation of $v$ as advanced time), so both families have $\d r < 0$
--- all future-directed null rays move to smaller $r$.  The horizon is
a one-way membrane: light can enter but not escape.
\physnote{The tilting of light cones across the horizon is the geometric
content of the event horizon.  No local physics changes at $r = 2M$;
the causal structure simply closes off escape routes to infinity.}

%
\begin{figure}[htbp]
  \centering
  % ef-light-cones.tex — Light cones in Schwarzschild and ingoing EF coordinates
% Inputable TikZ fragment for fig:ef-light-cones
% Requires: tikz with calc, arrows.meta (loaded by ehbook.cls)
%
% Two-panel comparison showing how light cones tilt at the event horizon.
% Panel (a): Schwarzschild (r, t) — symmetric cones pinching at r=2M.
% Panel (b): Ingoing EF (r, v)  — ingoing edge horizontal, outgoing tilts inward.
%
% Physics (M=1): null directions from ds²=0.
%   (a) dt/dr = ±1/(1−2/r).  Symmetric about time axis.
%   (b) Ingoing: dv=0.  Outgoing: dv/dr = 2/(1−2/r).
%
% Scales: s_r = 0.7 cm/M, s_t = s_v = 0.35 cm/M.  Arm length L = 0.35 cm.
%
\begin{tikzpicture}[
    >={Stealth[length=5pt]},
    font=\footnotesize,
  ]

  % ── Panel (a): Schwarzschild coordinates ───────────────────────────

  \begin{scope}[shift={(0,0)}]

    % Axis lengths
    \def\rmin{0.5}    % left edge in cm (r ≈ 0.7M)
    \def\rmax{3.5}    % right edge in cm (r ≈ 5M)
    \def\tmax{2.6}    % top of t axis

    % Horizon position: r=2M → x = 2×0.7 = 1.4 cm
    \def\rh{1.4}

    % Gray shaded strip for r < 2M (coordinates undefined)
    \fill[EHMarginGray, opacity=0.06]
      (0, -0.3) rectangle (\rh, \tmax);

    % Horizon line
    \draw[EHHorizonCrimson, dashed, thick]
      (\rh, -0.3) -- (\rh, \tmax)
      node[above, font=\scriptsize, text=EHHorizonCrimson] {$r = 2M$};

    % Axes
    \draw[->, EHMarginGray, thin] (-0.15, 0) -- (\rmax, 0)
      node[right, font=\scriptsize] {$r$};
    \draw[->, EHMarginGray, thin] (0, -0.15) -- (0, \tmax)
      node[above, font=\scriptsize] {$t$};

    % ── Light cones ──────────────────────────────────────────────────
    % Each cone: filled triangle with base at (r, 0), edges ±(Δx, Δy).
    % Pre-computed from plan: dt/dr = ±1/f, f = 1−2/r, arm length 0.35 cm.

    % r = 4.2M → x = 2.94
    \fill[EHAccretionGold, opacity=0.15]
      (2.94, 0) -- ++(0.253, 0.242) -- ++(-0.506, 0) -- cycle;
    \draw[EHAccretionGold, thick]
      (2.94, 0) -- ++(0.253, 0.242)  (2.94, 0) -- ++(-0.253, 0.242);
    \fill[EHAccretionGold] (2.94, 0) circle (1.2pt);

    % r = 3.2M → x = 2.24
    \fill[EHAccretionGold, opacity=0.15]
      (2.24, 0) -- ++(0.210, 0.280) -- ++(-0.420, 0) -- cycle;
    \draw[EHAccretionGold, thick]
      (2.24, 0) -- ++(0.210, 0.280)  (2.24, 0) -- ++(-0.210, 0.280);
    \fill[EHAccretionGold] (2.24, 0) circle (1.2pt);

    % r = 2.6M → x = 1.82
    \fill[EHAccretionGold, opacity=0.15]
      (1.82, 0) -- ++(0.147, 0.318) -- ++(-0.294, 0) -- cycle;
    \draw[EHAccretionGold, thick]
      (1.82, 0) -- ++(0.147, 0.318)  (1.82, 0) -- ++(-0.147, 0.318);
    \fill[EHAccretionGold] (1.82, 0) circle (1.2pt);

    % r = 2.2M → x = 1.54
    \fill[EHAccretionGold, opacity=0.15]
      (1.54, 0) -- ++(0.063, 0.344) -- ++(-0.126, 0) -- cycle;
    \draw[EHAccretionGold, thick]
      (1.54, 0) -- ++(0.063, 0.344)  (1.54, 0) -- ++(-0.063, 0.344);
    \fill[EHAccretionGold] (1.54, 0) circle (1.2pt);

    % Panel label
    \node[font=\scriptsize, anchor=north] at (1.75, -0.55) {(a) Schwarzschild};

  \end{scope}

  % ── Panel (b): Ingoing Eddington–Finkelstein coordinates ──────────

  \begin{scope}[shift={(5.0, 0)}]

    % Axis lengths
    \def\rmin{0.2}
    \def\rmax{3.5}
    \def\vmax{2.6}

    % Horizon position: r=2M → x = 1.4 cm
    \def\rh{1.4}

    % Horizon line
    \draw[EHHorizonCrimson, dashed, thick]
      (\rh, -0.3) -- (\rh, \vmax)
      node[above, font=\scriptsize, text=EHHorizonCrimson] {$r = 2M$};

    % Axes
    \draw[->, EHMarginGray, thin] (-0.15, 0) -- (\rmax, 0)
      node[right, font=\scriptsize] {$r$};
    \draw[->, EHMarginGray, thin] (0, -0.15) -- (0, \vmax)
      node[above, font=\scriptsize] {$v$};

    % ── Light cones ──────────────────────────────────────────────────
    % Ingoing edge: dv=0, always (−0.35, 0) — horizontal pointing left.
    % Outgoing edge: dv/dr = 2/(1−2/r), pre-computed per plan.

    % r = 4.2M → x = 2.94
    \fill[EHAccretionGold, opacity=0.15]
      (2.94, 0) -- ++(-0.350, 0) -- ++(0.513, 0.310) -- cycle;
    \draw[EHAccretionGold, thick]
      (2.94, 0) -- ++(-0.350, 0)  (2.94, 0) -- ++(0.163, 0.310);
    \fill[EHAccretionGold] (2.94, 0) circle (1.2pt);

    % r = 3.2M → x = 2.24
    \fill[EHAccretionGold, opacity=0.15]
      (2.24, 0) -- ++(-0.350, 0) -- ++(0.473, 0.328) -- cycle;
    \draw[EHAccretionGold, thick]
      (2.24, 0) -- ++(-0.350, 0)  (2.24, 0) -- ++(0.123, 0.328);
    \fill[EHAccretionGold] (2.24, 0) circle (1.2pt);

    % r = 2.6M → x = 1.82
    \fill[EHAccretionGold, opacity=0.15]
      (1.82, 0) -- ++(-0.350, 0) -- ++(0.429, 0.341) -- cycle;
    \draw[EHAccretionGold, thick]
      (1.82, 0) -- ++(-0.350, 0)  (1.82, 0) -- ++(0.079, 0.341);
    \fill[EHAccretionGold] (1.82, 0) circle (1.2pt);

    % r = 2.0M → x = 1.40 (at the horizon — outgoing is vertical)
    \fill[EHAccretionGold, opacity=0.15]
      (1.40, 0) -- ++(-0.350, 0) -- ++(0.350, 0.350) -- cycle;
    \draw[EHAccretionGold, thick]
      (1.40, 0) -- ++(-0.350, 0)  (1.40, 0) -- ++(0, 0.350);
    \fill[EHAccretionGold] (1.40, 0) circle (1.2pt);

    % r = 1.5M → x = 1.05 (inside — outgoing tilts inward)
    \fill[EHAccretionGold, opacity=0.15]
      (1.05, 0) -- ++(-0.350, 0) -- ++(0.239, 0.332) -- cycle;
    \draw[EHAccretionGold, thick]
      (1.05, 0) -- ++(-0.350, 0)  (1.05, 0) -- ++(-0.111, 0.332);
    \fill[EHAccretionGold] (1.05, 0) circle (1.2pt);

    % r = 1.0M → x = 0.70 (deep inside — both edges point left)
    \fill[EHAccretionGold, opacity=0.15]
      (0.70, 0) -- ++(-0.350, 0) -- ++(0.103, 0.247) -- cycle;
    \draw[EHAccretionGold, thick]
      (0.70, 0) -- ++(-0.350, 0)  (0.70, 0) -- ++(-0.247, 0.247);
    \fill[EHAccretionGold] (0.70, 0) circle (1.2pt);

    % Annotation: arrow showing both edges point to smaller r
    \draw[EHMarginGray, thin, ->]
      (0.25, 0.85) -- ++(-0.20, 0)
      node[left, font=\scriptsize, text=EHMarginGray, align=right]
      {both edges\\[-1pt]point inward};

    % Panel label
    \node[font=\scriptsize, anchor=north] at (1.75, -0.55) {(b) Eddington--Finkelstein};

  \end{scope}

\end{tikzpicture}

  \caption[Light cones in Schwarzschild and Eddington--Finkelstein coordinates]{%
    Light cones in Schwarzschild coordinates (left) and ingoing
    Eddington--Finkelstein coordinates (right) for the Schwarzschild geometry
    with $M = 1$.  In Schwarzschild coordinates, the light cones pinch shut
    at $r = 2M$ as $g_{tt} \to 0$ and $g_{rr} \to \infty$, and the
    coordinates do not extend through the horizon (shaded strip).  In
    Eddington--Finkelstein coordinates, the ingoing null direction
    $\d v = 0$ maintains a constant slope, while the outgoing direction
    tilts progressively inward; at $r = 2M$, the outgoing cone is tangent
    to the horizon.  Inside, both null directions point to smaller $r$,
    showing that the horizon is a one-way causal boundary.}
  \label{fig:ef-light-cones}
\end{figure}

\paragraph{Outgoing Eddington--Finkelstein coordinates.}

To complete the picture, an analogous construction with the retarded
time $u = t - r_*$ yields outgoing Eddington--Finkelstein coordinates.
The substitution $\d t = \d u + \d r_*$ produces a cross-term with the
opposite sign:
%
\begin{equation}
  \d s^2 = -\!\left(1 - \frac{2M}{r}\right)\d u^2
           - 2\,\d u\,\d r
           + r^2\,\d\Omega^2 \,.
  \label{eq:ef-outgoing}
\end{equation}
%
This metric covers a different patch of the maximally extended
spacetime: the exterior and the white hole interior (but not the
black hole interior or the second asymptotic region).  For the ray
tracer, outgoing coordinates are less useful than ingoing ones ---
backward-traced photons typically fall inward toward the black hole ---
but they appear naturally in the Kruskal--Szekeres construction
(\cref{sec:schw-kruskal}).
\xrefnote{The complete causal structure, including both ingoing and
outgoing patches, is developed in the Penrose diagram of
\cref{sec:schw-kruskal}.}

\paragraph{Why not Eddington--Finkelstein?}

These coordinates solve the horizon problem, so why doesn't the
codebase use them?  Two reasons.  First, the
metric (\cref{eq:ef-ingoing}) is written in spherical coordinates
$(v, r, \theta, \varphi)$, which introduce their own coordinate
singularity at $\theta = 0$ and $\theta = \pi$ --- the geodesic
integrator would need special treatment at the poles.  Second, the
angular coordinates break the flat-spacetime limit: at large $r$ the
metric approaches $\eta_{\mu\nu}$ in spherical form, not in the
Cartesian form that the codebase uses to set up initial conditions.

The codebase instead uses \emph{Kerr--Schild coordinates}, which
express the Schwarzschild metric in Cartesian spatial coordinates
$(t, x, y, z)$ as a flat background plus a single null-vector
correction (\cref{eq:kerr-schild-form}).  This representation is
horizon-penetrating, pole-free, and extends naturally to the
spinning Kerr metric --- the subject of the next section.
The Kerr--Schild form is less standard in the pedagogical literature,
so some geometric intuition is harder to read off directly ---
the Eddington--Finkelstein picture developed above remains the
clearest way to see why the horizon is a one-way membrane.
\implnote{The Haskell type \texttt{SchwarzschildKS} stores only the
mass $M$ and constructs all metric components from the Kerr--Schild
decomposition.  The polymorphic metric function
\texttt{schwarzschildMetricFn} works with any \texttt{Floating} type,
enabling automatic differentiation for Christoffel computation without
hand-coded derivatives.}
